\section{Método e critérios da revisão}
Foi realizada análise documental dos arquivos fornecidos: repositório do projeto, entrega final da Atividade XII, orientação oficial da XIII e versões Alfa A--D, com atenção especial às versões mais recentes. A revisão bibliográfica reuniu artigos científicos revisados por pares, trabalhos de economia de plataformas, estudos de experimentação empreendedora e fontes institucionais do setor pet e de normalização acadêmica. A estrutura, as citações e as referências seguem o guia institucional da PUC Minas, 6ª edição, articulado às normas ABNT vigentes para trabalhos acadêmicos, citações e referências (NBR 14724:2024, NBR 10520:2023 e NBR 6023:2025), além das normas de resumo, sumário e numeração progressiva (NBR 6028:2021, NBR 6027:2012 e NBR 6024:2012) \cite{pucmanual2025}.

Os critérios de seleção foram: (i) pertinência direta a plataformas digitais, crescimento de startups, experimentação e economia unitária; (ii) identificação bibliográfica rastreável por periódico, DOI ou instituição; (iii) atualidade quando a informação é conjuntural; e (iv) distinção entre evidência acadêmica, dado setorial e hipótese produzida pela equipe. A pesquisa não constitui revisão sistemática com protocolo PRISMA; trata-se de revisão narrativa orientada à decisão do projeto.
